\scanpage{341}
\mspage{8}

dans une catégorie $\deleted{\unsure{X}}\unsure{Y|V}$, telle que la \deleted{\unsure{somme}} $\varinjlim$ $L=b\unsure{\amalg|+}_a c$ \uline{dans $\deleted{\unsure{X}}Y^\wedge$} soit un objet non asphérique de $\deleted{\unsure{X}}Y^\wedge$, i.e. $\deleted{\unsure{X}}Y_{/L}$ soit non asphérique.

\deleted{\unsure{Je|J'en}} \unsure{sens|sais} qu'il faut, pour un tel contre-exemple, que $a\unsure{\mathbin{-}|\longrightarrow}b$ et $a\mathbin{-}c$ ne soient pas des mono. L'idée, c'est de prendre un exemple universel où ces flèches ne sont pas mono, en prenant

\deleted{\diagram{$Y=\Biggl($\begin{tikzcd}[ampersand replacement=\&,row sep=tiny,execute at end picture={\draw[line width=.35pt] ([xshift=-12pt,yshift=-4pt]C.south west) -- ([xshift=4pt,yshift=4pt]B.north east);}]
 \& |[alias=B]|\deleted{b}\\
 a\arrow[ur,bend left=23,"\unsure{\beta'_0|\beta_0}"]\arrow[ur,bend right=17,"\beta_1"']\arrow[dr,bend left=16,"\gamma_1"]\arrow[dr,bend right=24,"\gamma_0"'] \&\\
 \& |[alias=C]|c
\end{tikzcd}$\Biggr)$}}

\deleted{catégorie avec trois objets $a,b,c$, \unsure{les|ses} flèches non identiques étant $\beta_0,\beta_1,\gamma_0,\gamma_1$. On trouve que $L\in\mathrm{Ob}\,Y^\wedge$ est donné par}

\deleted{\[
 \begin{aligned}
 L(b)&\simeq L(c)=e\quad\text{(ens. à 1 élément)}\\
 L(a)&=\underbrace{\{\bar\beta_0,\bar\beta_1,\bar\gamma_0,\bar\gamma_1\}}_{\substack{\text{\unsure{ens en corr.}}\\\text{\unsure{biun.} avec }\{\beta_0,\beta_1,\gamma_0,\gamma_1\}}}\big/\ \text{relation d'équivalence}\\
 &\hspace{7em}\text{qui identifie }\beta_0\text{ et }\gamma_0\text{ en un \unsure{$=$|même} élément }\unsure{\varphi|\psi|\phi}_0\\
 &=\{\deleted{\unsure{\bar\gamma}}\varphi_0,\deleted{\unsure{\bar\gamma_1|\bar\beta_1}}\bar\beta_1,\bar\gamma_1\}
 \end{aligned}
\]}

\deleted{\[
 \begin{aligned}
 \text{et}\qquad L(\beta_i)&=\bar\beta_i,\quad L(\gamma_i)=\bar\gamma_i\qquad\text{pour }i\in\{0,1\},\\
 \text{donc}\qquad L(\beta_0)&=L(\gamma_0)=\varphi_0\\
 L(\beta_1)&=\bar\beta_{\unsure{1|2}},\quad L(\gamma_1)=\bar\gamma_1.
 \end{aligned}
\]}

\deleted{La catégorie $Y_{/L}$ au dessus de $Y$ a donc comme fibres en $b,c,a$ la cat. finale (deux fois) et la catégorie à trois éléments $\varphi_0,\bar\beta_1,\unsure{\bar\beta_1|\bar\gamma_1}$.}

% archive sheet number: 341
